\section{Đặt D-Tracker trong bối cảnh và đọc phản biện}
\label{sec:discussion}

\subsection{Related work trong một trang}

Paper chia công trình liên quan thành ba nhóm. Hiểu ba nhóm này giúp thấy D-Tracker kế thừa gì và
cố giải quyết điều gì.

\paragraph{Tensor analysis.}
Tucker, PARAFAC và PARAFAC2 là các cách phân rã tensor quen thuộc; nhiều biến thể online đã được
phát triển cho dữ liệu lớn. MLDS ghép linear dynamical system với Tucker; SMF và SSMF tập trung
vào mùa vụ và dự báo. DISMO ghép PARAFAC với phương trình vi phân phi tuyến, cho kết quả diễn giải
tốt, nhưng \emph{không} mô hình hóa lan truyền giữa các địa điểm. D-Tracker thêm đúng phần còn
thiếu đó bằng reaction--diffusion.

\paragraph{Time-series forecasting.}
Các phương pháp thống kê cổ điển gồm autoregression và Kalman filter. Họ mô hình sâu dùng CNN, RNN,
Transformer như PatchTST và Autoformer; LaST/CoST học biểu diễn trend/mùa vụ. Ưu điểm là độ chính
xác trên nhiều bài toán, nhược điểm là khó diễn giải và thường cần huấn luyện lại tốn kém, không
phù hợp trực tiếp với stream.

\paragraph{Social-activity modeling.}
Các mô hình web search từng nghiên cứu quan hệ cạnh tranh giữa từ khóa. FluxCube dùng ý tưởng
physics-informed neural network để học diffusion trong tensor, nhưng tốn thời gian vì phải ước lượng
nhóm địa điểm lặp lại. D-Tracker đặt diffusion vào một hệ phương trình nhỏ, cập nhật tăng dần.

\begin{table}[htbp]
  \centering
  \scriptsize
  \caption{So sánh khả năng theo paper Table~1. ``Có'' là đáp ứng, ``--'' là không được báo cáo
  là đáp ứng. D-Tracker là phương pháp duy nhất đáp ứng cả sáu yêu cầu theo bảng này.}
  \label{tab:capabilities}
  \begin{tabular}{@{}lccccccc@{}}
    \toprule
    & DISMO & FluxCube & PatchTST & Autoformer & DeepAR & LaST/CoST & D-Tracker \\
    \midrule
    Forecasting & Có & Có & Có & Có & Có & Có & Có \\
    Seasonality & Có & Có & -- & Có & -- & Có & Có \\
    Diffusion & -- & Có & -- & -- & -- & -- & Có \\
    Independent of length & -- & -- & -- & -- & -- & -- & Có \\
    Parameter free & Có & -- & -- & -- & -- & -- & Có \\
    Outlier & -- & -- & -- & -- & -- & -- & Có \\
    \bottomrule
  \end{tabular}
\end{table}

\subsection{Điểm mạnh thật sự}

Điểm mạnh không chỉ là ``dự báo chính xác hơn''. D-Tracker kết hợp ba thứ thường tách rời:
\begin{enumerate}
  \item một mô hình \textbf{sinh} có tham số diễn giải được (trend tăng/giảm, diffusion giữa nhóm);
  \item một cơ chế \textbf{streaming} chỉ dùng cửa sổ gần nhất;
  \item một tiêu chuẩn \textbf{MDL} để chọn cấu trúc thay vì thử tay.
\end{enumerate}
Với người làm phân tích, lợi ích là có thể nhìn một chuỗi dài và hỏi ``dynamics nào đang chi phối,
ở nhóm nào, lan theo hướng nào?'' thay vì chỉ nhận một con số dự báo.

\subsection{Giả định và giới hạn cần đọc kèm}

Đây là phần đánh giá của tài liệu, không phải các kết luận được paper chứng minh:

\begin{description}[leftmargin=1.8em,style=nextline]
  \item[Diễn giải không phải nhân quả.] Hệ reaction--diffusion là một \emph{giả định cấu trúc}. Nó
  có thể khớp một pattern trễ do yếu tố thứ ba gây ra. Paper không có thiết kế nhân quả, biến công
  cụ hay thí nghiệm can thiệp.
  \item[Phép cộng ba thành phần.] Giả định $\mathcal{X}^{c}\approx\widehat{\mathcal{X}}_d+
  \widehat{\mathcal{X}}_s+\widehat{\mathcal{X}}_o$ là mạnh. Nếu trend và mùa vụ tương tác phi tuyến,
  hoặc outlier không thưa, mô hình
  có thể phải dùng outlier quá nhiều và mất tính diễn giải.
  \item[Hệ số không âm.] Ràng buộc $\mathbf{W}\ge0$ giúp đọc nhóm và giữ dấu, nhưng cũng hạn chế
  các quan hệ cần trọng số âm (ví dụ một từ khóa làm giảm quan tâm từ khóa khác) trừ khi chúng
  được biểu diễn gián tiếp qua nhiều dynamics.
  \item[Không duy nhất.] Phân rã tensor thường không xác định duy nhất; các nhóm có thể xoay/hoán vị
  mà vẫn cho cùng tensor xấp xỉ. Vì vậy, đừng quá tập trung vào nhãn ``Group 1''; hãy đọc cấu trúc
  tương đối và kiểm tra tính ổn định.
  \item[Cửa sổ cố định.] Một cửa sổ $L_c$ cố định có thể bỏ lỡ chu kỳ dài hơn cửa sổ hoặc một sự
  kiện xảy ra trước đó. Chọn $L_c$ quá ngắn hoặc quá dài đều có hậu quả; paper dùng 104 tuần cho
  GoogleTrends và 56 ngày cho Covid-19.
  \item[Không có bất định.] Mô hình trả về điểm dự báo, không có khoảng tin cậy hay phân phối dự báo.
  Với quyết định rủi ro cao, cần bổ sung đánh giá bất định.
  \item[Dữ liệu thiếu và không đều.] Paper không trình bày cơ chế chính thức cho missing values,
  timestamp không đều hoặc thay đổi định nghĩa địa điểm. Đây là khoảng trống cần xử lý khi áp dụng.
\end{description}

\subsection{Độ tin cậy của bằng chứng thực nghiệm}

Các kết quả có điểm mạnh: sáu stream thật, bảy baseline, nhiều horizon, ablation cho từng module và
phân tích thời gian chạy. Nhưng cần đọc kèm:
\begin{itemize}
  \item không có kiểm định thống kê hay khoảng tin cậy cho chênh lệch MAE;
  \item ``tối đa 57\%'' không nêu rõ mẫu số tại chỗ, nên không thể diễn giải như một effect size
        thống nhất;
  \item tốc độ 8.0 lần gắn với phần cứng và cách huấn luyện baseline cụ thể;
  \item ablation cho thấy một số module không cải thiện ở một số dataset, đúng như paper thừa nhận;
  \item kết quả Covid-19 chỉ có hai tín hiệu (ca nhiễm, tử vong) và không đại diện cho toàn bộ động
        học dịch tễ.
\end{itemize}
Những điều này không phủ định kết quả; chúng giới hạn phạm vi kết luận.

\subsection{Khi nào nên và không nên dùng D-Tracker?}

\textbf{Nên cân nhắc} khi dữ liệu là tensor nhiều địa điểm/từ khóa theo thời gian, cần dự báo ngắn
hạn, cần gợi ý về nhóm và hướng lan truyền, và muốn cập nhật khi dữ liệu mới đến.

\textbf{Không nên dùng một mình} khi mục tiêu là suy luận nhân quả, khi dữ liệu thiếu/không đều
nghiêm trọng, khi cần khoảng tin cậy, hoặc khi số chiều quá nhỏ đến mức mô hình tensor không có lợi
thế. Trong những trường hợp đó, D-Tracker có thể là một \emph{nguồn giả thuyết} hữu ích, không phải
câu trả lời cuối cùng.

\subsection{Những câu hỏi phản biện nên chuẩn bị}

\begin{enumerate}
  \item Nếu hai vùng tăng cùng lúc do một chiến dịch toàn cầu, mô hình có luôn gán là diffusion không?
  \item Khi số regime tăng, $\mathcal{F}$ có thể lớn đến mức nào, và chi phí bộ nhớ ra sao?
  \item Các nhóm địa điểm có ổn định qua các lần chạy/khởi tạo khác nhau không?
  \item Điều gì xảy ra nếu chọn $L_c$ ngắn hơn chu kỳ mùa vụ?
  \item Vì sao dùng MDL thay vì cross-validation cho cấu trúc?
  \item Kết quả có còn đúng nếu missing values được điền theo cách khác?
\end{enumerate}
Đây là các câu hỏi hợp lệ để thảo luận, không phải lỗi chắc chắn của paper. Một số câu paper trả lời
một phần; số khác cần thực nghiệm bổ sung.

\begin{keypoint}
Giá trị của D-Tracker nằm ở \textbf{cách nó kết nối mô hình sinh, nén MDL và streaming}, không phải
ở một con số MAE đơn lẻ. Đọc paper công bằng nghĩa là vừa ghi nhận kết quả mạnh, vừa giữ rõ giả định
và giới hạn của chúng.
\end{keypoint}

\begin{selfcheck}
\begin{enumerate}
  \item Trong paper Table~1, khả năng nào chỉ D-Tracker có?
  \item Vì sao ``diffusion không cải thiện ở Language'' lại là thông tin hữu ích, không phải lỗi cần che?
  \item Nêu hai giả định có thể bị vi phạm khi áp dụng D-Tracker vào dữ liệu thực tế.
  \item Nếu cần kết luận nhân quả, cần thêm loại bằng chứng nào ngoài mô hình này?
\end{enumerate}
\end{selfcheck}
